\chapter{Participant Questionnaire}
\label{app:questionnaire}

This appendix reproduces the final participant-facing questionnaire used during live data collection.

\begingroup
\small
\setlength{\parindent}{0pt}
\setlength{\parskip}{0.5em}
\setlength{\tabcolsep}{4pt}

\newcolumntype{Y}{>{\raggedright\arraybackslash}X}
\newcommand{\formbox}{\fbox{\rule{0pt}{1.1ex}\rule{1.1ex}{0pt}}}
\newcommand{\formline}{\underline{\hspace{3cm}}}
\newcommand{\formanswerline}{\par\vspace{2.5em}\hrulefill\par}
\newcommand{\question}[1]{\par\medskip\textbf{#1}\par}
\newcommand{\responselegend}{\textit{Response options: SD = Strongly Disagree, D = Disagree, N = Neutral, A = Agree, SA = Strongly Agree.}}

\textbf{Participant ID:} \formline \hfill \textbf{Date:} \formline

\section*{Screening and Background}

\question{1. Are you 18 or older?}
\formbox\ Yes \hspace{2em} \formbox\ No

\question{2. Have you previously studied computer architecture or computer organisation?}
\formbox\ Yes \hspace{2em} \formbox\ No

\question{3. Have you used a VR headset before?}
\formbox\ Yes \hspace{2em} \formbox\ No

\question{4. How experienced are you with Virtual Reality?}
\formbox\ None \hspace{1em}
\formbox\ Limited \hspace{1em}
\formbox\ Moderate \hspace{1em}
\formbox\ Good \hspace{1em}
\formbox\ High

\question{5. Before this session, how confident are you in your understanding of computer architecture concepts?}
\formbox\ Extremely confident \hspace{0.6em}
\formbox\ Somewhat confident \hspace{0.6em}
\formbox\ Neutral \hspace{0.6em}
\formbox\ Somewhat not confident \hspace{0.6em}
\formbox\ Extremely not confident

\question{6. How difficult do you usually find topics such as instruction flow, pipelining, stack behaviour, and memory hierarchy?}
\formbox\ Extremely difficult \hspace{0.6em}
\formbox\ Somewhat difficult \hspace{0.6em}
\formbox\ Neutral \hspace{0.6em}
\formbox\ Somewhat easy \hspace{0.6em}
\formbox\ Extremely easy

\question{7. Which of the following topics have you previously studied?}
\textit{Select all that apply.}

\begin{tabularx}{\linewidth}{YYY}
\formbox\ Instruction Flow & \formbox\ Pipelining & \formbox\ Stack \\
\formbox\ Memory Hierarchy & \formbox\ Assembly & \formbox\ None of the above \\
\end{tabularx}

\section*{System Experience}

\question{8. Please indicate how strongly you agree with the following statements for System Usability.}
\responselegend

\begin{tabular}{p{0.50\linewidth}ccccc}
 & \textbf{SD} & \textbf{D} & \textbf{N} & \textbf{A} & \textbf{SA} \\
The VR system was easy to learn. & \formbox & \formbox & \formbox & \formbox & \formbox \\
I could navigate the environment without difficulty. & \formbox & \formbox & \formbox & \formbox & \formbox \\
The controls felt intuitive. & \formbox & \formbox & \formbox & \formbox & \formbox \\
The visual presentation made the concepts easier to follow. & \formbox & \formbox & \formbox & \formbox & \formbox \\
I could focus on the learning task without being distracted by the interface. & \formbox & \formbox & \formbox & \formbox & \formbox \\
I would feel confident using this system again. & \formbox & \formbox & \formbox & \formbox & \formbox \\
I experienced discomfort, motion sickness, or visual fatigue while using the system. & \formbox & \formbox & \formbox & \formbox & \formbox \\
\end{tabular}

\question{9. Please indicate how strongly you agree with the following statements for System Engagement.}
\responselegend

\begin{tabular}{p{0.50\linewidth}ccccc}
 & \textbf{SD} & \textbf{D} & \textbf{N} & \textbf{A} & \textbf{SA} \\
The experience held my attention. & \formbox & \formbox & \formbox & \formbox & \formbox \\
I felt actively involved in the learning process. & \formbox & \formbox & \formbox & \formbox & \formbox \\
The system made the topic feel more interesting. & \formbox & \formbox & \formbox & \formbox & \formbox \\
I was motivated to explore the concepts further. & \formbox & \formbox & \formbox & \formbox & \formbox \\
The interactive format improved my focus compared with standard learning materials. & \formbox & \formbox & \formbox & \formbox & \formbox \\
\end{tabular}

\question{10. Please indicate how strongly you agree with the following statements about your learning experience.}
\responselegend

\begin{tabular}{p{0.50\linewidth}ccccc}
 & \textbf{SD} & \textbf{D} & \textbf{N} & \textbf{A} & \textbf{SA} \\
The VR experience helped me understand the topic better. & \formbox & \formbox & \formbox & \formbox & \formbox \\
The visualisations helped me build a mental model of what was happening. & \formbox & \formbox & \formbox & \formbox & \formbox \\
The system made abstract processes easier to understand. & \formbox & \formbox & \formbox & \formbox & \formbox \\
I would prefer this as a supplement to traditional teaching materials. & \formbox & \formbox & \formbox & \formbox & \formbox \\
I believe I learned something useful from this session. & \formbox & \formbox & \formbox & \formbox & \formbox \\
\end{tabular}

\section*{Knowledge Check}

Answer the following questions to assess your understanding of the computer architecture concepts covered in the session.

\textbf{Reference Guide}

\begin{itemize}
\item MIPS: Microprocessor without Interlocked Pipeline Stages, a Reduced Instruction Set Computer (RISC) architecture commonly used for teaching computer architecture.
\item \texttt{add}: Adds the values from two registers and stores the result in a destination register.
\item \texttt{addi}: Adds an immediate (constant) value to a register and stores the result.
\item \texttt{lw} (Load Word): Loads a 32-bit word from memory into a register.
\item \texttt{sw} (Store Word): Stores a 32-bit word from a register into memory.
\item \texttt{beq} (Branch if Equal): Branches to a target address if two register values are equal.
\item \texttt{j} (Jump): Unconditionally jumps to a specified instruction address.
\end{itemize}

\question{11. Which register stores the address of the next instruction?}
\begin{tabularx}{\linewidth}{YY}
\formbox\ Stack Pointer & \formbox\ Program Counter \\
\formbox\ Instruction Register & \formbox\ Base Register \\
\end{tabularx}

\question{12. Which stage of the instruction cycle reads an instruction from memory?}
\begin{tabularx}{\linewidth}{YY}
\formbox\ Decode & \formbox\ Fetch \\
\formbox\ Execute & \formbox\ Write-back \\
\end{tabularx}

\question{13. During instruction execution, the instruction register primarily stores:}
\begin{tabularx}{\linewidth}{YY}
\formbox\ Program variables & \formbox\ The currently executing instruction \\
\formbox\ The stack pointer value & \formbox\ Cache addresses \\
\end{tabularx}

\question{14. What happens to the program counter after an instruction is fetched?}
\begin{tabularx}{\linewidth}{YY}
\formbox\ It resets to zero & \formbox\ It increments to the next instruction address \\
\formbox\ It copies into the stack pointer & \formbox\ It stores the result of the ALU \\
\end{tabularx}

\question{15. Which pipeline stage usually performs arithmetic operations?}
\begin{tabularx}{\linewidth}{YY}
\formbox\ Fetch & \formbox\ Decode \\
\formbox\ Execute & \formbox\ Write-back \\
\end{tabularx}

\question{16. Branch instructions primarily affect which component?}
\begin{tabularx}{\linewidth}{YY}
\formbox\ Cache memory & \formbox\ Program counter \\
\formbox\ Stack pointer & \formbox\ Instruction register \\
\end{tabularx}

\question{17. What does this instruction do?}
\texttt{add \$t0, \$t1, \$t2}\\
\formbox\ Adds \texttt{\$t1} and \texttt{\$t2} and stores the result in \texttt{\$t0}

\formbox\ Adds \texttt{\$t0} and \texttt{\$t1} and stores the result in \texttt{\$t2}

\formbox\ Moves \texttt{\$t1} into \texttt{\$t2}

\formbox\ Multiplies registers

\question{18. What does this instruction do?}
\texttt{addi \$t0, \$zero, 10}\\
\formbox\ Stores 10 into \texttt{\$t0}

\formbox\ Clears \texttt{\$t0}

\formbox\ Moves \texttt{\$t0} into \texttt{\$zero}

\formbox\ Adds two registers

\question{19. Which instruction type loads data from memory into a register?}
\begin{tabularx}{\linewidth}{YY}
\formbox\ ADD & \formbox\ STORE \\
\formbox\ LOAD & \formbox\ JUMP \\
\end{tabularx}

\question{20. What is the role of register \$zero in MIPS?}
\begin{tabularx}{\linewidth}{YY}
\formbox\ Stores the program counter & \formbox\ Always contains the value 0 \\
\formbox\ Stores return addresses & \formbox\ Tracks the stack pointer \\
\end{tabularx}

\question{21. What does this instruction do?}
\texttt{lw \$t0, 0(\$s0)}\\
\formbox\ Stores \texttt{\$t0} into memory

\formbox\ Loads a word from memory into \texttt{\$t0}

\formbox\ Moves \texttt{\$s0} into \texttt{\$t0}

\formbox\ Clears memory

\question{22. What does this instruction do?}
\texttt{sw \$t0, 4(\$s0)}\\
\formbox\ Value from \texttt{\$t0} is stored in memory

\formbox\ Value from memory is loaded into \texttt{\$t0}

\formbox\ \texttt{\$s0} is overwritten

\formbox\ Stack pointer moves

\question{23. What does this instruction do?}
\texttt{beq \$t0, \$t1, label}\\
\formbox\ Jump if \texttt{\$t0 > \$t1}

\formbox\ Jump if \texttt{\$t0 == \$t1}

\formbox\ Jump if \texttt{\$t0 < \$t1}

\formbox\ Always jump

\question{24. What does this instruction do?}
\texttt{j label}\\
\formbox\ Jumps to \texttt{label} unconditionally

\formbox\ Jumps only if two registers are equal

\formbox\ Stores the return address in \texttt{\$ra}

\formbox\ Loads a value from memory

\question{25. Put the instruction stages in the correct order.}
\begin{tabularx}{\linewidth}{YY}
\formbox\ Execute, Fetch, Decode, Write-back & \formbox\ Fetch, Decode, Execute, Write-back \\
\formbox\ Decode, Fetch, Execute, Write-back & \formbox\ Fetch, Execute, Decode, Write-back \\
\end{tabularx}

\question{26. How confident are you in your answers to the previous section?}
\formbox\ Not at all confident \hspace{1em}
\formbox\ 1 \hspace{1em}
\formbox\ 2 \hspace{1em}
\formbox\ 3 \hspace{1em}
\formbox\ 4 \hspace{1em}
\formbox\ 5 \hspace{1em}
Very confident

\section*{Open-Ended Feedback}

\question{27. What part of the experience helped your understanding most?}
\formanswerline

\question{28. What part was confusing or ineffective?}
\formanswerline

\question{29. What would you improve in the system?}
\formanswerline

\question{30. Were there any moments where the VR interface got in the way of learning?}
\formanswerline

\question{31. Would you use this again as a learning aid? Why or why not?}
\formanswerline

\endgroup
